%% guia-exemplo.tex
%% Documento mínimo de teste — Guia Visual CRP-SP
%% Compilar com: lualatex guia-exemplo.tex
%%
%% Em produção (TeX Live 2024+), adicionar antes de \documentclass:
%%   \DocumentMetadata{
%%     lang = pt-BR, pdfversion = 2.0,
%%     pdfstandard = ua-2, testphase = {phase-III},
%%   }

\documentclass[guia]{crpsp_acessivel}

% --- Paleta desta publicação ---
\crpspCorPrincipal{8B1A4A}
\crpspCorSecundaria{2E5090}
\crpspCorAcessoria{D4A843}

\title{Guia de Exemplo CRP-SP}
\author{Conselho Regional de Psicologia de São Paulo}
\date{2026}

\begin{document}

% ──────────────────────────────
\frontmatter
% ──────────────────────────────

\begin{titlepage}
\centering
\vspace*{3cm}
{\fontsize{24}{28}\selectfont\bfseries\color{cor1}
Guia de Exemplo\\[0.5em]
CRP-SP\par}
\vfill
{\large 2026\par}
\end{titlepage}

% ──────────────────────────────
\mainmatter
% ──────────────────────────────

\chapter{Introdução ao Guia}

Este é o primeiro capítulo do guia de exemplo. O texto começa nesta
página, após a abertura dedicada do capítulo.

\section{O que é este guia}

Este guia foi produzido pelo CRP-SP como material de orientação
profissional. Ele segue os padrões de acessibilidade PDF/UA-2.

\begin{itemize}
  \item Primeiro item da lista com bullet colorido.
  \item Segundo item da lista.
  \begin{itemize}
    \item Sub-item de primeiro nível.
    \item Outro sub-item.
  \end{itemize}
  \item Terceiro item.
\end{itemize}

\section{Tabela de exemplo}

\begin{crpsptabela}{lll}
  Coluna A & Coluna B & Coluna C \\
  Dado 1   & Dado 2   & Dado 3   \\
  Dado 4   & Dado 5   & Dado 6   \\
  Dado 7   & Dado 8   & Dado 9   \\
\end{crpsptabela}

\chapter{Orientações Práticas}

Conteúdo do segundo capítulo. Aqui o redator pode usar o carimbo
de seção antes de uma nova seção temática:

\carimbosecao

\section{Procedimentos recomendados}

Texto corrido com as orientações práticas do guia.

\begin{enumerate}
  \item Primeira orientação.
  \item Segunda orientação.
  \item Terceira orientação.
\end{enumerate}

% ──────────────────────────────
\backmatter
% ──────────────────────────────

\chapter*{Referências}

Referências bibliográficas e normativas do guia.

\end{document}
